\documentclass[10pt]{article}
\usepackage[margin=20mm]{geometry}
\usepackage{amsmath,amssymb,mathtools,fontspec,unicode-math,xurl,hyperref}
\IfFontExistsTF{DejaVu Serif}{\setmainfont{DejaVu Serif}}{\setmainfont{Latin Modern Roman}}
\setmathfont{latinmodern-math.otf}
\setlength{\parindent}{0pt}\setlength{\parskip}{6pt}\sloppy
\title{Pointwise kernel values and atomic energies:
a technical verification report}
\author{AI-assisted technical review}
\date{}
\begin{document}
\maketitle
Abstract: This report examines the acquired author version of Shu and Wang’s generalized Erdős–Turán inequalities. An explicit bounded, even, lower-semicontinuous kernel satisfies every stated hypothesis of the energy stability theorem but has zero energy at a Dirac mass whose Wasserstein-infinity distance from the uniform measure is \(1/2\). A related construction disproves the stated \(L^\infty\) perturbation conclusion for arbitrary probability measures. These are two manifestations of one substantive issue: almost-everywhere kernel data do not determine energies of atomic measures. A separate local endpoint estimate in the multiplier proof requires a logarithm; the stated multiplier bound survives the repair. No counterexample to the principal Riesz-potential distance inequality is established. Scientific adjudication remains pending.\par
\section*{Identity, version, and inspected scope}
The reviewed paper is Ruiwen Shu and Jiuya Wang, Generalized Erdős–Turán inequalities and stability of energy minimizers, arXiv:2110.03019v1, submitted 6 October 2021, 23 pages [paper]. Its associated journal record is Journal of the European Mathematical Society, DOI 10.4171/JEMS/1583, published online on 3 January 2025 [journal]. This report concerns the exact acquired v1 PDF, not an established textual match to the subscription-restricted journal version. All page references below are one-based PDF pages; their printed numbers agree. The separately retained author-hosted PDF has not been substituted for this source.\par
The complete v1 text was read. Checks traced the transport characterization (Theorem 1.7), weighted Hall argument, regularized test-function construction, Fourier multiplier estimates, one-dimensional endpoint regimes, sharpness construction, and the energy applications. Particular attention was given to the quantification over all probability measures in \(\mathcal M(\mathbb T^d)\). External cited proofs and the particle simulation underlying Figure 1 were not independently reproduced. No article-specific executable code was found in the acquired source search; its absence is not a defect.\par
\section*{Definitions and the affected proof chain}
Equation reference: eq:energy, eq:h3, eq:spectral. The torus has normalized Lebesgue measure and its usual quotient metric. For an actual pointwise kernel \(W\), the paper defines \[
 \mathcal E_W[\rho]=\frac12\iint_{\mathbb T^d\times\mathbb T^d}
 W(x-y)\,d\rho(x)\,d\rho(y).\] The hypotheses on p. 4 require \(W\) to be integrable, even, lower semicontinuous and bounded below (H1), with zero mean and strictly positive nonzero Fourier coefficients (H2), and \[
 \frac1{|B(x;r)|}\int_{B(x;r)}(C_1+W(y))\,dy
 \le C(C_1+W(x))\] for some constants and every ball (H3). The additional lower bound is \(\widehat W(k)\ge c|k|^{-d+s}\), \(k\ne0\), with \(s<d\). Theorem 1.4 derives energy stability from the asserted identity \[
 \mathcal E_W[\rho]=\frac12\sum_{k\ne0}\widehat W(k)|\widehat\rho(k)|^2,\] which is equation (1.12) on p. 3. For \(d=1,s<0\), the conclusion is \(d_\infty(\rho,1)\lesssim \mathcal E_W[\rho]^{1/(2-s)}\). Theorem 1.6 on p. 5 then asserts the same control for any minimizer of \(\mathcal E_{\widetilde W}\), replacing the energy by \(\|W-\widetilde W\|_{L^\infty}\).\par
\section*{Substantive findings}
Problem 1 (Almost-everywhere kernel control does not control atomic energy).  The stated hypotheses do not justify [eq:spectral] (\#eq:spectral) for arbitrary probability measures. Consequently Theorem 1.4(ii)(a) and the stated perturbation theorem admit explicit counterexamples in the reviewed version.\par
\subsection*{A counterexample satisfying every energy-theorem hypothesis}
Equation reference: eq:w. Take \(d=1,s=-1\), and let \[K(x)=2\pi^2\left(x^2-x+\frac16\right),\qquad 0\le x\le1,\] with periodic extension. It is continuous and even on \(\mathbb T\). Direct integration gives zero mean, and two integrations by parts give \(\widehat K(k)=k^{-2}\) for every nonzero integer \(k\). Moreover \[-\frac{\pi^2}{6}\le K(x)\le\frac{\pi^2}{3},
 \qquad K(0)=\frac{\pi^2}{3},\quad K(1/2)=-\frac{\pi^2}{6}.\] Define the actual pointwise function \[
 W(x)=K(x)\quad(x\ne0\text{ in }\mathbb T),\qquad W(0)=0.\] This bounded even function is lower semicontinuous: at zero its assigned value is below its limiting value, and elsewhere it is continuous. Thus (H1) holds. Changing one point leaves all Lebesgue Fourier coefficients unchanged, proving (H2) and the quantitative lower bound with \(c=1\).\par
For a fully explicit check of (H3), set \(C_1=1+\pi^2/6\) and \(C=1+\pi^2/2\). Pointwise, \[1\le C_1+W(x)\le C.\] Every ball average in [eq:h3] (\#eq:h3) is therefore at most \(C\), whereas its right-hand side is at least \(C\). This verifies (H3) for all centers and radii, including the modified point. No unlisted regularity condition is needed.\par
Now take \(\rho=\delta_0\), which is permitted by the definition of \(\mathcal M\) on p. 8. Equation [eq:energy] (\#eq:energy) gives \[\mathcal E_W[\delta_0]=\frac12W(0)=0,\qquad
 d_\infty(\delta_0,1)=\frac12.\] The distance follows because the only coupling with first marginal \(\delta_0\) has the uniform second marginal, whose support reaches distance \(1/2\) from zero. Theorem 1.4(ii)(a) would require \[\frac12\le C\,0^{1/3}=0.\] The contradiction is independent of the implicit constant. The spectral expression instead equals \[\frac12\sum_{k\ne0}k^{-2}=\frac{\pi^2}{6},\] so the precise failed bridge is [eq:spectral] (\#eq:spectral). The claim that the uniform measure minimizes this energy also fails: for \(\nu=(\delta_0+\delta_{1/2})/2\), \[\mathcal E_W[\nu]=\frac14W(1/2)=-\frac{\pi^2}{24}<0=\mathcal E_W[1].\]\par
\subsection*{An independent failure of the perturbation estimate}
Keep the unperturbed kernel equal to the continuous \(K\), and set \[\widetilde W(x)=K(x)\ (x\ne0),\qquad \widetilde W(0)=-\pi^2.\] This is again bounded, even and lower semicontinuous. In the usual Lebesgue meaning of \(L^\infty\), \(\|K-\widetilde W\|_{L^\infty}=0\). Since \(\widetilde W\ge-\pi^2\) pointwise, \(\mathcal E_{\widetilde W}[\rho]\ge-\pi^2/2\) for every probability measure, with equality at \(\delta_0\). Thus the measure in this counterexample is a genuine global minimizer, not an arbitrary competing measure. Theorem 1.6 would again bound its distance \(1/2\) by zero.\par
The uniform-measure comparison (1.18) remains valid. The decisive error is the use of the Lebesgue essential supremum in (1.19) against \(d\rho(x)d\rho(y)\): for the Dirac measure, the difference \(K(0)-\widetilde W(0)=4\pi^2/3\) is fully sampled. This failure persists even though the unperturbed kernel is canonical and continuous.\par
\subsection*{Alternative readings, impact, and verified repair boundaries}
No global convention in the read v1 text restricts \(\mathcal M\) to densities, forces general \(W\) to be continuous, or specifies a canonical pointwise representative. Atomic measures are explicitly used in Section 2. Reading \(L^\infty\) as a pointwise supremum would remove the second counterexample but does not remove [eq:w] (\#eq:w), which already invalidates Theorem 1.4 without any perturbation norm.\par
This supports Major impact on the advertised energy applications: it disproves their stated guarantees within their quantified class. The main distance inequality in Theorem 1.1 and the canonical Riesz-kernel calculations do not depend on this invalid general energy identification. The examples do not challenge the numerical clustering pictures.\par
A precise restricted repair is available. If \(\rho=f\,dx\) with \(f\ge0\), \(\int f=1\), \(f\in L^2(\mathbb T^d)\), then \(W\in L^1\) gives absolute integrability of [eq:energy] (\#eq:energy), by \(\||W|*f\|_2\le\|W\|_1\|f\|_2\). Parseval then proves [eq:spectral] (\#eq:spectral); the original application of Theorem 1.1 yields Theorem 1.4 on this smaller measure class. If the perturbed minimizer is in this class and the uniform density remains an admissible competitor, the essential-supremum energy comparison is also valid. This proves a narrowed result, not the original assertion for atoms.\par
Alternatively, for a base kernel whose energy lower bound has been justified on all measures, imposing \(\widetilde W=W+V\) with \(|V(x)|\le\delta\) at every point gives \(|\mathcal E_{\widetilde W}[\rho]-\mathcal E_W[\rho]|\le\delta/2\). The minimizer comparison then gives \(\mathcal E_W[\rho]\le\delta\). This repairs the perturbation step only; the base kernel condition remains essential. No full repair for every singular kernel allowed by the original hypotheses is claimed.\par
Problem 2 (A missing logarithm in an intermediate multiplier estimate).  For the allowed parameter \(\beta=1\), the radial integral used in the last inequality of (3.26), p. 14, is not bounded by the displayed expression uniformly in \(|x|\).\par
Writing \(R=|x|>1\), the relevant nonnegative scalar integral is exactly \[\int_{R^{-1}<|\xi|\le1}|\xi|^{\beta-d-1}\,d\xi
 =|\mathbb S^{d-1}|\int_{R^{-1}}^1r^{\beta-2}\,dr.\] At \(\beta=1\) it equals \(|\mathbb S^{d-1}|\log R\), while the asserted bound is a constant multiple of \(R^{-(\beta-1)}+1=2\). This is a local invalid inequality in the proof, rather than a counterexample to Lemma 3.6.\par
Retaining the logarithm yields \(|u(x)|\le C R^{-d-1}(1+\log R)\) at this endpoint. It is still integrable at infinity. The periodization argument (3.27)–(3.28) therefore retains its \(L^1\) bound. The \(L^\infty\) bound follows directly from the absolutely summable Fourier coefficients and Schwartz decay: \[\sum_{k\in\mathbb Z^d}|\widehat\psi(\epsilon k)|\,|k|
 \le C\epsilon^{-d-1}.\] Interpolation recovers exactly the bound of Lemma 3.6 at \(\beta=1\). The displayed decay estimate may admit a sharper proof, but the particular scalar-integral inequality in (3.26) cannot. The supported impact is Minor: the local derivation requires correction and the downstream multiplier estimate is restored without changing its statement.\par
\section*{Other checked reasoning and remaining scope}
The weighted transport argument accommodates irrational weights; a rational-weight approximation in \(d_\infty\) was not silently substituted. In Section 3 the crucial exponent is \[\frac{d-1}{d}-\frac1q=\frac1p-\frac1d\ge0
 \quad(1\le p\le d),\] and the regularized set and its complement contain balls at the relevant scale, yielding the stated \(r^{d+d/p-s}\) power. The one-dimensional Hilbert-transform step uses \(1<q<\infty\); the paper does not invoke its false \(L^1\) or \(L^\infty\) endpoint boundedness. The sharpness densities have zero perturbation mass and remain nonnegative; their moment cancellations give the asserted near- and far-field powers. These checks support the reported scope, not a formal verification of every cited theorem.\par
The current arXiv landing page inspected during this review lists only v1. The EMS landing page lists the journal identity and dates and requires a subscription for its PDF. No correction to the present issue was found in that bounded metadata check. Absence from these pages is not proof that no correction exists; the journal-final applicability of the findings remains unresolved.\par
\section*{Conclusion}
Item | Assessment | Repair verification | Adoption / remaining work\par
Problem 1 | Major in acquired v1 | Restricted density and pointwise-perturbation statements proved above | Author adoption unknown; compare journal-final text\par
Problem 2 | Minor proof estimate | Logarithmic endpoint repair preserves Lemma 3.6 | Author adoption unknown\par
Adjudication | Pending | Initial AI checks only | No human review supplied\par
9 R. Shu and J. Wang, Generalized Erdős–Turán inequalities and stability of energy minimizers, arXiv:2110.03019v1 (6 October 2021). https://arxiv.org/abs/2110.03019v1. EMS Press, article metadata for the same title, DOI 10.4171/JEMS/1583 (https://doi.org/10.4171/JEMS/1583), published online 3 January 2025; landing page consulted 26 October 2026. https://ems.press/journals/jems/articles/14298463.\par
\section*{Supporting checks}
The counterexamples and repairs are analytic calculations, not numerical tests or executions of author software.\par
`acea4137e7182300ed00323101bbcbd88bdf9e9e3edd8af704590bbca8486599`.\par
The controller’s UI observation is separately attributed; backend identity was not independently verified. These are initial AI checks, not independent acceptance or human adjudication.\par
\end{document}
